%=====================================================================
\section{Discussion}
\label{sec:discussion}

\paragraph{What the static network already tells us.}
Conclusions that survive aggregation can be read from \cref{tab:cross}.
For CIT, the final aggregate describes degree heterogeneity, clustering and
connected components. It cannot by itself identify preferential attachment or
ageing: these compare new citations with prior degrees and paper ages.
Even densification requires a sequence of historical $N(t),E(t)$ pairs, rather
than one final graph. The PS/CM claims remain to be tested.

\paragraph{What requires time.}
Some conclusions need windows (Parts I and III) and others the full event order
(Parts II and IV).
\expect{[Expected] Burstiness, turnover, community reorganisation, time-respecting
reachability and spreading speed cannot be obtained from the aggregate.}

\paragraph{Why network types differ.}
The three networks differ in generating mechanism (physical proximity, voluntary
messaging, citing earlier work), time scale and edge semantics (a contact ends, a
message is a point event, a citation persists).
\expect{[Expected] Proximity is constrained by timetables and space and thus has
strong periodic structure, messages are driven by individual bursty decisions.}
CIT has persistent edges and measured age-dependent target preferences.
A cross-network ordering still needs PS/CM measurements.

\paragraph{Robustness.}
We check the choice of window size, the observation period (first vs.\ second
half), the community-detection resolution and the choice of null model.
For CIT, the full-period densification exponent is 1.343 but falls to 1.065 in
the later half, while remaining 1.427 with cutoff 2015. Positive degree preference
persists under all declared variants, although its exponent depends on the degree
offset. The diameter verdict is unchanged by the sampled-source sensitivity
checks. These are qualitative robustness findings, not evidence that coverage
and censoring have been removed.


% Owner: B (measured citation-specific interpretation).
\paragraph{Citation-specific evidence.}
CIT supports degree preference and age-dependent target choice, but separates two
claims often grouped together: cumulative edges grow superlinearly, while the
sampled effective diameter does not shrink. A three-parameter citation curve
fits selected histories well, yet its paired MAE advantage over linear prediction
is unresolved. Annual zero clustering is imposed by the strict-year snapshot
construction. These examples show why representation, fitted-sample description
and held-out prediction require separate interpretations. Late reference coverage,
within-venue selection and right censoring constrain all citation conclusions.

\paragraph{Limitations.}
There is one dataset per network type, so generality is limited and conclusions
hold ``in these data''. CIT also has within-venue sampling, annual timestamp ties and changing reference
coverage. Its year-block bootstrap intervals are descriptive; cumulative graph
observations are dependent. The 20-draw null ensembles give coarse tail resolution.
Forecast selection excludes low-cited papers, and 38.6\% of short-history WSB fits
hit parameter bounds; fitted collapse is not independent validation.
Boundary effects
and sampling bias remain (finite two-day contact window, sensor participation).
Citation networks are degenerate temporal networks (DAG, no deletion), which is
why several methods do not apply. Each system also contains several time scales
\citep{holme2015modern}, so a window that suits one process may hide another. Finally, SocioPatterns data are released under
CC BY-NC-SA.
